\documentclass[11pt]{article}
\usepackage[a4paper,margin=1in]{geometry}
\usepackage[T1]{fontenc}
\usepackage{lmodern}
\usepackage{csquotes}
\usepackage[backend=biber,style=authoryear,sorting=nyt,maxbibnames=99]{biblatex}

\addbibresource{annotated-bibliography.bib}

\title{Annotated Bibliography\\\large Generated from Existing Knowledge Base}
\author{}
\date{\today}

\begin{document}
\maketitle

This annotated bibliography is generated from the current repository knowledge base.
Each entry includes a standard citation and a short annotation describing why the source is relevant.

\section*{Short Summary of Each Paper}
\begin{itemize}
  \item \textcite{abicht2023owl}: Reviews the current practical usability and maintenance status of OWL reasoners.
  \item \textcite{acpo2011gpgde}: Defines the ACPO principles and foundational UK good-practice guidance for handling digital evidence.
  \item \textcite{akinbi2022blockchainiot}: Synthesizes blockchain-based IoT forensic investigation models and their adoption constraints.
  \item \textcite{alharbi2011proactivereactive}: Proposes a model that combines proactive evidence collection with reactive investigation.
  \item \textcite{barahona2006reasoningweb}: Explains core Semantic Web reasoning approaches and rules-ontology integration patterns.
  \item \textcite{bratsas2024ctikg}: Surveys how knowledge graphs and Semantic Web tooling are used in cyber threat intelligence.
  \item \textcite{casey2015cybox}: Introduces CybOX-based standardization for digital forensic information exchange and provenance.
  \item \textcite{casey2017case}: Presents CASE as a community specification to improve cyber-investigation interoperability.
  \item \textcite{casino2022reviewofreviews}: Consolidates major trends, gaps, and future priorities across digital forensics review literature.
  \item \textcite{coppolillo2024llasp}: Evaluates fine-tuning strategies that improve LLM performance on answer set programming tasks.
  \item \textcite{cybok2021v110}: Provides the CyBOK taxonomy used for structuring cybersecurity and forensics knowledge areas.
  \item \textcite{elsappagh2015cbr}: Reviews case representation methods in case-based reasoning, including ontology-driven approaches.
  \item \textcite{erdem2016aspapps}: Summarizes practical ASP applications and discusses benefits and open challenges.
  \item \textcite{fahndrich2023strongai}: Assesses how strong-AI concepts are framed and evidenced in digital forensics literature.
  \item \textcite{gebser2016clingo5}: Describes clingo 5 architecture for integrating ASP with theory-specific reasoning.
  \item \textcite{goldstrawwhite2022legalpolicy}: Maps legal and policy obligations affecting UK digital forensics practice.
  \item \textcite{hall2022xai}: Frames explainable AI levels and deployment implications for forensic investigations.
  \item \textcite{javed2022comprehensivesurvey}: Offers a broad survey of computer forensics methods, tools, and future directions.
  \item \textcite{fsr2015g218}: Outlines a lifecycle approach for validating digital forensic methods in regulated settings.
  \item \textcite{liang2025kag}: Introduces KAG for combining structured knowledge and LLMs in professional-domain QA.
  \item \textcite{lifschitz2008whatisasp}: Provides a foundational introduction to the principles and modeling style of ASP.
  \item \textcite{marcopoulos2018onlinesparc}: Presents onlineSPARC, a browser-based environment for ASP development and learning.
  \item \textcite{kadage2024aienhanceddf}: Describes AI-enabled automation opportunities for faster forensic investigation workflows.
  \item \textcite{parsonage2025agafa}: Proposes the AGAFA neuro-symbolic framework for explainable forensic analysis.
  \item \textcite{pietranik2020ontologymappings}: Defines metrics for evaluating ontology mappings at concept and instance levels.
  \item \textcite{sikos2021ontologydf}: Reviews ontology engineering as an enabler for AI-supported digital forensics.
  \item \textcite{silva2024ontologiesdf}: Reviews digital-forensics ontology research and provides recommendations for practice.
  \item \textcite{syed2016uco}: Introduces UCO to unify heterogeneous cybersecurity knowledge representations.
  \item \textcite{tageldin2023mlforensics}: Surveys machine-learning methods used in smart-environment forensic investigations.
  \item \textcite{tok2025scopeontology}: Extends UCO/CASE with smart-city ontology concepts for threats and forensics.
  \item \textcite{tully2020qualitydf}: Analyzes quality and accreditation impacts on digital forensics in England and Wales.
\end{itemize}

\printbibliography[title={References with Annotations}]

\end{document}
